\documentclass[dvipdfmx]{article} % 文章の形式を設定
\usepackage[margin=2.5cm]{geometry} % 書式の空白を設定
\usepackage[utf8]{inputenc} % 文字コードをUTF-8に設定
\usepackage{hyperref} % 目次にリンクを付けるため
\usepackage{lipsum} % ダミーテキスト用
\usepackage{tcolorbox} % 枠を利用するため
\usepackage{amsmath} % 数式の記述を行うため
\usepackage{bm} % ベクトルを太字で表示するため
\usepackage{graphicx} % 画像を表示するため
\usepackage{float} % 画像正しい位置で表示するため
\usepackage{tensor} % テンソルを記載するため
\usepackage{multicol} % 複数段落を作成するため
\usepackage{tikz} % 図を作成するため

\title{シュバルツシルト時空における光源の軌道予測}
\author{20041054　大豆生田 幹}
\date{}
\begin{document}
\maketitle 
概要\\
\noindent
ブラックホールに関する観測は近年積極的に行われており、主だったニュースとして2019年4月、国際協力ブロジェクト「イベント・ホライズン・テレスコープ（Event Horizon Telescope）」にて、ブラックホールの影が撮影されたことが挙げられる。\\

\noindent
このような話題に興味があったので、私は「ブラックホールに落ちる天体がどのように見えるかを予測すること」を研究の目標とした。\\
これを行うには\\
\indent
$(1)$ ブラックホールの周りを回る、光源である天体がどのような軌道を描くのか\\
\indent
$(2)$ 天体から放たれた光は重力の影響を受け、どのような軌道で観測者に届くのか\\
を考える必要がある。\\

\noindent
本発表は $(1)$ のについて\\
つまり「ブラックホールの周りを回る、光原である天体がどのような軌道を描くのか」を予想する。\\

\noindent
「天体の軌道予測」というと、解析力学で学習したようなケプラー問題のようなものを想像されるかもしれないが、そのように予測した計算結果を観測結果を照らすと、誤差が出てしまうことが知られている。\\
そこで、今回は一般相対論を用いて厳密に軌道を予測することを試みる。\\

\noindent
ただし、ブラックホールの周りを回る天体の軌道といっても、宇宙にはさまざまな質量、性質をもつ物質があり、それを考慮するといくらでも複雑な軌道が考えられてしまう。\\
そこで今回は、曲率が時間変化せず、ブラックホールを中心として球対称である、という状況を仮定した元で導出される「シュバルツシルト時空」と呼ばれる時空での天体の動きを計算した。\\

具体的に、以下の流れで発表を行う。\\
\indent
$1.$ 重い星の周り動く星の軌道をニュートン力学を用いて考える\\
\indent
$2.$ 重い星の周り動く星の軌道を一般相対論を用いて考える\\
\indent
$3.$ 数値解析を用いて予測した運動を見てみる\\

参考文献\\
$(1)$　佐々木 節　「一般相対論」　産業図書　2011（第7版）\\
$(2)$　Robert M. Wald　「General Relativity」　The University of Chicago　1984

\end{document}